\section{K-017: a sharp-period bound cannot be deduced from freeness alone}

This section records a failed route, not a semistable counterexample and
not a resolution of an all-dimensional main question.  The attempted
extension was to bound the two-primary determinant character by
$2^{n-1}$ for every free tensor support, without using semistability.
The following exact example refutes that stronger statement.

\begin{proposition}[K-017]
There is a free support in $\{0,\ldots,7\}^3$ whose diagonal
stabilizer contains an element of order256 with determinant-product
character of order256.  Its complete fixed tensor space is unstable
for $SL_8^3$.
\end{proposition}

\begin{proof}
Let $\zeta$ be a primitive256th root, and use the three exponent lists
\[
\begin{aligned}
A&=(198,34,133,229,237,252,232,5),\\
B&=(251,222,34,4,65,19,42,0),\\
C&=(29,24,241,218,249,58,157,0).
\end{aligned}
\]
Let $g$ be the triple of diagonal matrices with these eigenvalues.
The total exponent is2933, congruent to117 modulo256, so both
$g$ and $\chi(g)$ have order256.  The complete zero-weight support is
\[
\begin{gathered}
021,044,075,117,133,146,216,245,324,343,362,\\
401,422,457,537,552,563,600,615,671,707,710.
\end{gathered}
\]
Here a three-digit word denotes its three coordinate indices.  This
list is obtained by testing $A_i+B_j+C_k\equiv0\pmod{256}$ for all
512 triples; the retained integer verifier repeats that finite test.
All entries in each exponent list are distinct.  Therefore two
zero-weight edges agreeing in two indices also agree in the third,
and this support is free.

For an explicit instability certificate put
\[
\begin{aligned}
y_A&=(14,-6,1,33,9,12,8,1),\\
y_B&=(-1,6,-6,-12,13,-9,-30,0),\\
y_C&=(-7,-8,-3,18,-27,-14,-7,0).
\end{aligned}
\]
The sum of their entries is $-15$.  On the listed support,
$y_{A,i}+y_{B,j}+y_{C,k}$ equals64 at343 and zero at every other
edge.  Thus a nonnegative margin-one solution would give the
contradiction
\[
  -15=\sum_{i,j,k}q_{ijk}(y_{A,i}+y_{B,j}+y_{C,k})\ge0.
\]
Equivalently, replacing each factor list $y_J$ by
$8y_J-(\sum y_J){\bf1}$ gives integral traceless lists.  Their
sum on each support edge is15, except at343 where it is527.
The associated one-parameter subgroup of $SL_8^3$ therefore sends
every tensor in the complete fixed space to zero.  This proves
instability of the whole fixed space, rather than only infeasibility
of one selected basis.
\end{proof}

The incidence columns also have Smith form $(1^{21},256)$, checked
exactly and retained with the support.  This additional computation
is not needed for the preceding counterexample: the displayed
element and its primitive character already suffice.

\paragraph{Other attempted routes and their precise stopping points.}
At $n=5$, write $\pi=1-\zeta_{16}$ and normalize
$v_\pi(2)=8$.  The exact residue K-014 has valuation $-117$,
whereas the symmetric chain contribution has valuation $-99$.
The sum of the non-chain contributions also has valuation $-117$.
Thus the chain does not dominate this sum in the2-adic valuation.
The four weighted terms of smallest valuation are spectral types
3,4,5,6, each of valuation $-123$; their lowest terms cancel.
The integer/Fraction computation is retained as
\texttt{residue\_2adic.py} and \texttt{residue\_2adic.json}.
These are finite cancellation diagnostics, with no extrapolation
to arbitrary dimension.

Ordinary positivity of $k_n(s)$ at attained integer degrees is
also insufficient for the definition of $s(k_n)$ in GCT VI,
Definition1.2.4, printed p.~14: that definition tests every
nonzero constituent polynomial at every positive integer argument.
In particular the ordinary-positivity theorem of Amanov and
Yeliussizov (arXiv:2202.11059v2, Theorem8.12, printed p.~31),
which applies to even $s\ge2(n+1)^2$, does not bound the
off-congruence-class roots of those constituent polynomials.
This statement of the distinction was checked directly against
both sources in this run.
